% from ega4-2-fr.tex
\oldpage[IV-2]{180}
où $X=\operatorname{Spec}(A)$ (avec $A=\mathcal O_{X,x}$) est un schéma local.
Notons que $A'=\mathcal O_{X',x'}$ est par hypothèse un $A$-module
fidèlement plat, donc contenant $A$, et par suite l'hypothèse que $X$ est
géométriquement unibranche au point $x$, ou l'hypothèse que $X'$ est
géométriquement unibranche au point $x'$, entraînent toutes deux que $A$
est un anneau local intègre (puisque $A$, étant réduit, est aussi isomorphe
à un sous-anneau de $A'_{\mathrm{red}}$). Soient $K$ le corps des fractions
de $A$, $B$ la clôture intégrale de $A$, et posons
$Y=\operatorname{Spec}(B)$; il revient au même de dire que $X$ est
géométriquement unibranche au point $x$ ou que le morphisme $g:Y\to X$
est radiciel au point $x$ (6.15.3). Soit
$Y'=Y\otimes_k k'=Y\times_X X'$ de sorte que l'on a le diagramme
commutatif
\[
\xymatrix{
Y \ar[d]_g & Y' \ar[l]^h \ar[d]_{g'}\\
X & X' \ar[l]_f
}
\]
Notons que $f$ est un morphisme plat; donc (6.15.4.1) le morphisme entier
$g'_{\mathrm{red}}$ est birationnel. D'autre part, comme $Y$ est normal,
$Y'$ est géométriquement unibranche (6.15.6). Pour que $X'$ soit
géométriquement unibranche en $x'$, il faut et il suffit donc que $g'$
soit radiciel au point $x'$ (6.15.5). Or, $x$ est la projection de $x'$ dans
$X$ et il est équivalent de dire que $g$ est radiciel au point $x$ ou que
$g'$ est radiciel au point $x'$ (6.15.3.1); enfin, pour que $X$ soit
géométriquement unibranche au point $x$, il faut et il suffit que $g$ soit
radiciel en ce point, d'où la proposition.
\end{proof}

\begin{lemma}[6.15.8]
\label{IV.2.s6.15.8-fr}
Soient $k$ un corps séparablement clos (autrement dit, tel que la clôture
algébrique de $k$ soit radicielle sur $k$), $X$ un $k$-préschéma localement
de type fini sur $k$, $x$ un point fermé de $X$. Si $X$ est unibranche au
point $x$, il est géométriquement unibranche en ce point.
\end{lemma}

\begin{proof}
En effet, on sait (I, 6.4.2) que $k(x)$ est une extension algébrique de
$k$; comme le corps résiduel de la clôture intégrale de
$(\mathcal O_{X,x})_{\mathrm{red}}$ (clôture intégrale qui est par
hypothèse un anneau local) est une extension algébrique de $k(x)$, c'est
une extension radicielle de $k$ par hypothèse, donc aussi de $k(x)$.
\end{proof}

\begin{corollary}[6.15.9]
\label{IV.2.s6.15.9-fr}
Soient $k$ un corps, $X$ un $k$-préschéma localement de type fini sur $k$,
$k'$ une extension séparablement close de $k$ (autrement dit, telle que la
clôture algébrique de $k'$ soit radicielle sur $k'$); alors, pour que $X$
soit géométriquement unibranche, il faut et il suffit que
$X'=X\otimes_k k'$ soit unibranche.
\end{corollary}

\begin{proof}
Compte tenu de (6.15.7), on est ramené à prouver que si $X$ est unibranche
et $k$ séparablement clos, alors $X$ est géométriquement unibranche. Or, $X$
est géométriquement unibranche en ses points fermés (6.15.8); nous en
conclurons dans (10.4.9) que $X$ est géométriquement unibranche en tous ses
points, en nous appuyant sur le fait (signalé à la fin de (6.15.2)) que
l'ensemble des points où $X$ est géométriquement unibranche est
constructible (bien entendu, le cor. (6.15.9) ne sera pas utilisé
jusque-là).
\end{proof}

Ce résultat justifie, en une certaine mesure, la terminologie de
« géométriquement unibranche ».

